\section{Filtering by a Class}
\label{sec-filter}

\stitle{The hierarchy as a membership test}.
Collapse each cycle of $H$ into one class, pick for every class with several parents the
parent with the longest path from a root, and number the classes in postorder of the
resulting spanning forest.
In a forest the classes below $c$ are those whose number lies in one interval
$[\lo(c),\po(c)]$; in a DAG each class carries a \emph{set} of intervals $I(c)$, built
children first by merging its own interval with the sets of its children
(Agrawal et al.~\cite{agrawal1989}), and $|I(c)|\le 1+|E'(c)|$ for the $|E'(c)|$
non-tree edges below $c$.
The numbers are keys of the engine's B+ tree, $(\po(c),c,0)$ for each class, so
\begin{itemize}
\item the \emph{members} of $c$ are one range scan per interval of $I(c)$, in
  $O(|I(c)|)$ descents plus the members, and
\item whether a class $u$ lies below $c$ is a binary search of $\po(u)$ in $I(c)$, with
  no scan and no memory beyond the intervals.
\end{itemize}
On the four ontologies the mean class needs 1.3 to 2.0 intervals (at most 4,329, within the
bound), against 5.3 million ancestor pairs for ChEBI's closure and 402,305 intervals.
The labels are built in 4 to 69\,ms for 13,850 to 218,723 classes and rebuilt when the
hierarchy changes; term ids are untouched.

\stitle{Four strategies}.
A query $\mathrm{near}(q,c,k)$ can be answered in four ways, all returning terms of
$\desc(c)$ only.
\begin{description}
\item[Pre-filter] resolves $c$ to its members by range scans and scores them exactly:
  about $\gamma_m\,|\desc(c)|$ time, where $\gamma_m$ covers a scan step, a lookup from
  term to row, and one dot product.
\item[Post-filter] asks an index for $k'>k$ answers, keeps those below $c$ (a binary
  search each), and multiplies $k'$ by four until $k$ remain or the index is exhausted.
\item[Masked search] hands the index the predicate ``the row's number lies in an
  interval of $I(c)$'', so a list scan or a graph walk skips rows outside $c$ and
  nothing is resolved in advance; each row touched costs one binary search.
  Masked IVF keeps probing lists, nearest first, until it holds $k$ allowed rows or
  has probed all (Algorithm~\ref{alg-ivf}); masked HNSW still crosses rows it may not
  return, to keep the graph connected, and keeps a beam of $\ef$ allowed candidates.
\item[Auto] pre-filters when $|\desc(c)|\le\theta$ and otherwise runs masked HNSW.
\end{description}

\begin{algorithm}[t]
\caption{\textsc{Near}$(q,c,k)$: the auto rule}\label{alg-auto}
\small
\lIf{$c$ is not a class of $H$}{\Return $\emptyset$}
\eIf{$|\desc(c)|\le\theta$}{
  $R \gets$ the rows of the keys in the range scans $[(a,0,0),(b{+}1,0,0))$ for each $[a,b]\in I(c)$\;
  \Return the $k$ rows of $R$ with the largest $q\cdot\mathbf{v}$ \tcp*{pre-filter: exact}
}{
  \Return \textsc{MaskedHnsw}$(q,k,\ef,\;r\mapsto \po(r)\in I(c))$\;
}
\end{algorithm}

\begin{algorithm}[t]
\caption{\textsc{MaskedHnsw}$(q,k,\ef,\mathit{allowed})$}\label{alg-hnsw}
\small
descend the upper layers greedily from the entry point to the node nearest $q$ (the filter is not applied)\;
$\mathit{frontier} \gets \{$that node$\}$; $W \gets \emptyset$ \tcp*{$W$ holds allowed rows only}
\While{$\mathit{frontier}\ne\emptyset$}{
  $x \gets$ the frontier node nearest $q$\;
  \lIf{$|W|\ge\ef$ and $x$ is farther than the worst row of $W$}{\textbf{break}}
  \For{each unvisited neighbor $y$ of $x$ on layer 0}{
    add $y$ to $\mathit{frontier}$ whether or not it is allowed\;
    \lIf{$\mathit{allowed}(y)$}{add $y$ to $W$; drop the worst row if $|W|>\ef$}
  }
}
\Return the $k$ best rows of $W$\;
\end{algorithm}

\begin{algorithm}[t]
\caption{Masked IVF: the $k$ best rows allowed by the class}\label{alg-ivf}
\small
sort the $\nu$ lists by the similarity of their centroids to $q$\;
$\mathit{best} \gets \emptyset$; $\mathit{seen} \gets 0$\;
\For{each list $\ell$ in that order}{
  \For{each key $(\ell,r,0)$ of the prefix scan of $\ell$ in the B+ tree}{
    \lIf{the number of row $r$ lies in an interval of $I(c)$}{score $r$; keep it among the $k$ best}
  }
  $\mathit{seen} \gets \mathit{seen}+1$\;
  \lIf{$\mathit{seen}\ge\rho$ and $|\mathit{best}|=k$}{\textbf{break}}
}
\Return $\mathit{best}$\;
\end{algorithm}

\begin{lemma}[Exactness and completeness]\label{lem-exact}
Pre-filter returns exactly $\mathrm{near}(q,c,k)$. Masked IVF returns
$\min(k,|\desc(c)|)$ answers, and exactly $\mathrm{near}(q,c,k)$ when $\rho=\nu$.
\end{lemma}

\noindent
Pre-filter scores every member; masked IVF stops only with $k$ allowed rows or after
all lists, and with all lists probed it has scored every allowed row
(Appendix~\ref{apx-proofs}).
Masked HNSW (Algorithm~\ref{alg-hnsw}) crosses disallowed nodes to keep the graph
connected and keeps a beam of $\ef$ \emph{allowed} rows, so it does not stop early for
lack of allowed neighbors, which is also why it is slow for a small class: it walks until
it has found $\ef$ of them.
An index asked for $k$ answers and filtered afterwards has neither property.

\stitle{The filter is correlated with the vectors}.
Costing a filtered search by the filter's selectivity assumes independence: a filter
that passes a fraction $f(c)$ of the terms needs about $k/f(c)$ candidates.
In an ontology the names of a class's descendants resemble the class and one
another, so the neighbors of a term concentrate under the term's own ancestors.
Table~\ref{tab-correlation} measures it: for a query that is a member of $c$, or a synonym
of one, the share of the unfiltered top 100 under $c$ is 0.03 to 0.59, which is 10 to 100
times the fraction of terms below $c$ for classes of up to 1,000 members, and for a
query outside it falls to 0.00 to 0.24, below that fraction for the largest classes.
The estimate $k/f(c)$ is therefore too large for an inside query, which a post-filter
answers cheaply, and too small for an outside one, for which a post-filter's
candidates run out before $k$ allowed rows are seen.

\begin{table}[t]
\caption{Share of a query's unfiltered top 100 neighbors that lie under the class, by query kind and number of classes below $c$ (pooled mean; the independence estimate would be the fraction of terms below $c$).}
\label{tab-correlation}
\small
\begin{tabular}{@{}lrrrr@{}}
\toprule
Below $c$ & independent & inside & outside & synonym \\
\midrule
$\le$10 & 0.000 & 0.03 & 0.00 & 0.03 \\
11--100 & 0.001 & 0.14 & 0.00 & 0.11 \\
101--1,000 & 0.012 & 0.31 & 0.01 & 0.29 \\
1,001--5,000 & 0.054 & 0.54 & 0.03 & 0.43 \\
$>$5,000 & 0.334 & 0.59 & 0.24 & 0.46 \\
\bottomrule
\end{tabular}
\end{table}


\stitle{Resolve or test}.
Masked search can test the filter per row (one binary search over $I(c)$) or resolve it
first into a bitmap by range scans.
Resolving costs $\gamma_m|\desc(c)|$ whether or not the search ever looks at those
rows; testing costs only for rows the walk touches.
Pre-filtering resolves and scores, so it is the right plan exactly when
$\gamma_m|\desc(c)|$ is less than the masked walk's cost $\gamma_h$:
\[
\mathrm{pre\text{-}filter\ wins} \iff |\desc(c)| < \theta = \gamma_h/\gamma_m .
\]
The walk's cost depends on where the query is: for a query inside $c$ the walk finds
allowed rows at once, and for one outside it must cross rows it cannot return, so
$\gamma_h$ and $\theta$ differ, and a single threshold has to sit between them
(Section~\ref{sec-exp}).
